\documentclass[11pt,a4paper]{article}
% preamble.tex — estilo común de todos los apuntes Froth.
% Cada apunte hace % preamble.tex — estilo común de todos los apuntes Froth.
% Cada apunte hace % preamble.tex — estilo común de todos los apuntes Froth.
% Cada apunte hace \input{preamble} justo después de \documentclass.
\usepackage[margin=2.4cm]{geometry}
\usepackage{xcolor}
\definecolor{litblue}{RGB}{31,79,121}
\definecolor{litgreen}{RGB}{27,94,32}
\definecolor{litorange}{RGB}{230,81,0}
\usepackage{parskip}
\usepackage{titlesec}
\titleformat{\section}{\Large\bfseries\color{litblue}}{\thesection}{0.6em}{}
\titleformat{\subsection}{\large\bfseries\color{litblue}}{\thesubsection}{0.6em}{}
\usepackage{tcolorbox}
\usepackage{fancyhdr}
\pagestyle{fancy}
\fancyhf{}
\fancyhead[L]{\small\color{litblue}Froth — Apuntes de ML}
\fancyhead[R]{\small\thepage}
\renewcommand{\headrulewidth}{0.4pt}
\usepackage[colorlinks=true,linkcolor=litblue,urlcolor=litblue]{hyperref}

% Cabecera de cada apunte: \cabecera{numero}{titulo}
\newcommand{\cabecera}[2]{%
  \begin{tcolorbox}[colback=litblue,coltext=white,colframe=litblue,arc=2mm]
    {\Large\bfseries Apunte #1 — #2}\\[3pt]
    {\small Proyecto Froth · aprender Machine Learning construyendo}
  \end{tcolorbox}\vspace{0.8em}}

% Cajas temáticas
\newtcolorbox{concepto}[1]{colback=blue!4,colframe=litblue,title={Concepto: #1},arc=1.5mm}
\newtcolorbox{practica}{colback=green!5,colframe=litgreen,title={Pruébalo tú (teclado en mano)},arc=1.5mm}
\newtcolorbox{ojo}{colback=orange!8,colframe=litorange,title={Ojo / error típico},arc=1.5mm}
\newtcolorbox{resumen}{colback=gray!8,colframe=gray!60,title={En una frase},arc=1.5mm}

\newcommand{\codigo}[1]{\texttt{#1}}
 justo después de \documentclass.
\usepackage[margin=2.4cm]{geometry}
\usepackage{xcolor}
\definecolor{litblue}{RGB}{31,79,121}
\definecolor{litgreen}{RGB}{27,94,32}
\definecolor{litorange}{RGB}{230,81,0}
\usepackage{parskip}
\usepackage{titlesec}
\titleformat{\section}{\Large\bfseries\color{litblue}}{\thesection}{0.6em}{}
\titleformat{\subsection}{\large\bfseries\color{litblue}}{\thesubsection}{0.6em}{}
\usepackage{tcolorbox}
\usepackage{fancyhdr}
\pagestyle{fancy}
\fancyhf{}
\fancyhead[L]{\small\color{litblue}Froth — Apuntes de ML}
\fancyhead[R]{\small\thepage}
\renewcommand{\headrulewidth}{0.4pt}
\usepackage[colorlinks=true,linkcolor=litblue,urlcolor=litblue]{hyperref}

% Cabecera de cada apunte: \cabecera{numero}{titulo}
\newcommand{\cabecera}[2]{%
  \begin{tcolorbox}[colback=litblue,coltext=white,colframe=litblue,arc=2mm]
    {\Large\bfseries Apunte #1 — #2}\\[3pt]
    {\small Proyecto Froth · aprender Machine Learning construyendo}
  \end{tcolorbox}\vspace{0.8em}}

% Cajas temáticas
\newtcolorbox{concepto}[1]{colback=blue!4,colframe=litblue,title={Concepto: #1},arc=1.5mm}
\newtcolorbox{practica}{colback=green!5,colframe=litgreen,title={Pruébalo tú (teclado en mano)},arc=1.5mm}
\newtcolorbox{ojo}{colback=orange!8,colframe=litorange,title={Ojo / error típico},arc=1.5mm}
\newtcolorbox{resumen}{colback=gray!8,colframe=gray!60,title={En una frase},arc=1.5mm}

\newcommand{\codigo}[1]{\texttt{#1}}
 justo después de \documentclass.
\usepackage[margin=2.4cm]{geometry}
\usepackage{xcolor}
\definecolor{litblue}{RGB}{31,79,121}
\definecolor{litgreen}{RGB}{27,94,32}
\definecolor{litorange}{RGB}{230,81,0}
\usepackage{parskip}
\usepackage{titlesec}
\titleformat{\section}{\Large\bfseries\color{litblue}}{\thesection}{0.6em}{}
\titleformat{\subsection}{\large\bfseries\color{litblue}}{\thesubsection}{0.6em}{}
\usepackage{tcolorbox}
\usepackage{fancyhdr}
\pagestyle{fancy}
\fancyhf{}
\fancyhead[L]{\small\color{litblue}Froth — Apuntes de ML}
\fancyhead[R]{\small\thepage}
\renewcommand{\headrulewidth}{0.4pt}
\usepackage[colorlinks=true,linkcolor=litblue,urlcolor=litblue]{hyperref}

% Cabecera de cada apunte: \cabecera{numero}{titulo}
\newcommand{\cabecera}[2]{%
  \begin{tcolorbox}[colback=litblue,coltext=white,colframe=litblue,arc=2mm]
    {\Large\bfseries Apunte #1 — #2}\\[3pt]
    {\small Proyecto Froth · aprender Machine Learning construyendo}
  \end{tcolorbox}\vspace{0.8em}}

% Cajas temáticas
\newtcolorbox{concepto}[1]{colback=blue!4,colframe=litblue,title={Concepto: #1},arc=1.5mm}
\newtcolorbox{practica}{colback=green!5,colframe=litgreen,title={Pruébalo tú (teclado en mano)},arc=1.5mm}
\newtcolorbox{ojo}{colback=orange!8,colframe=litorange,title={Ojo / error típico},arc=1.5mm}
\newtcolorbox{resumen}{colback=gray!8,colframe=gray!60,title={En una frase},arc=1.5mm}

\newcommand{\codigo}[1]{\texttt{#1}}

\begin{document}
\cabecera{34}{Resumen extractivo vs abstractivo: por qué el nuestro no puede alucinar}

\section{El problema de la Fase 7}

El review v1 es un andamio honesto pero seco: keywords + listas de papers. Para que
\emph{diga} algo hay dos caminos, y la diferencia entre ellos es LA decisión de diseño
más importante de esta fase.

\begin{concepto}{extractivo vs abstractivo}
\textbf{Abstractivo}: un modelo GENERA texto nuevo que ``suena a'' resumen (ChatGPT,
Elicit). Fluido, pero el texto no existe en ninguna fuente: puede inventar afirmaciones
y --- peor --- referencias. \\
\textbf{Extractivo}: el sistema SELECCIONA frases reales que ya existen en los abstracts.
Menos fluido, pero cada frase tiene un paper de origen: la cita es trazable
\emph{por construcción}. Alucinar es imposible: no se genera ni una palabra.
\end{concepto}

\section{Por qué elegimos extractivo primero}

\begin{itemize}
  \item \textbf{Confianza}: un review académico con UNA referencia inventada muere.
        ChatGPT ``a pelo'' es famoso por citar papers que no existen.
  \item \textbf{Trazabilidad gratis}: como la frase viene de un abstract del corpus,
        el botón ``know more'' y la cita [n] ya funcionan sin verificación extra.
  \item \textbf{Es ML de verdad}: rankear frases usa embeddings, coseno, grafos y
        PageRank --- todo lo que ya construiste, a una escala nueva (frases, no papers).
  \item El pulido con LLM (parafrasear frases YA elegidas, verificando citas) queda como
        capa OPCIONAL encima --- nunca como base (7.6).
\end{itemize}

\section{La cocina: de abstract a frases utilizables}

En \texttt{froth/review.py}, \texttt{split\_sentences()} corta cada abstract donde hay
puntuación final seguida de mayúscula, y \texttt{\_self\_contained()} descarta lo que no
sobrevive solo: fragmentos ($<$8 palabras), maratones ($>$60) y frases que arrancan con
\textbf{anáforas} --- ``This method...'', ``It also...'' --- porque apuntan a un texto que
el lector del review no verá.

\begin{ojo}
El extractivo hereda las limitaciones de su materia prima: trabajamos con ABSTRACTS
(no full text), y un abstract mal escrito da frases mediocres. El review v2 sigue siendo
un andamio mejorado, no prosa final --- honestidad ante todo.
\end{ojo}

\begin{resumen}
Extractivo = seleccionar frases reales (trazable, sin alucinación); abstractivo = generar
texto (fluido, riesgoso). Froth construye sobre extractivo y deja el LLM local como
pulidor opcional de frases ya elegidas.
\end{resumen}

\end{document}
